% ==============================================================================
% Section: Institutional Background
% Status: NOT STARTED
% ==============================================================================

\section{Institutional Background}
\label{sec:institutional}

\TODO{Describe: PJM RPM structure, the Variable Resource Requirement (VRR) curve,
the Base Residual Auction (BRA) process, Locational Deliverability Areas (LDAs),
existing market power mitigation rules (Three Pivotal Supplier test, Capacity
Performance rules).}

% DATA NOTE: The panel has two structural features that must be acknowledged here:
%
% (1) AUCTION LEAD TIME DISRUPTION: PJM's BRA schedule was highly irregular across
%     the sample. Lead times ranged from ~13 months (2022/23, 2023/24) to effectively
%     zero (~1 month for 2024/25 — a catch-up auction held May 2024 for delivery
%     starting June 2024), then recovered to 10–27 months. No auction was held in
%     calendar year 2023. Discuss implications for forward capacity signals.
%
% (2) VRR CURVE REDESIGN (FERC ER25-1357): The VRR curve changed structurally
%     between the 2025/26 and 2026/27 delivery years. The old design had a $0 price
%     floor (3-point curve). The new design has an explicit price floor (~$177-179/
%     MW-day) and cap ($256.75/MW-day), converting to a 4-point curve with a flat
%     segment at the floor. The 2026/27 and 2027/28 BRAs both cleared at exactly the
%     VRR Point (a) price (the cap). This redesign is a natural experiment for the
%     VRR slope comparative static in Section~\ref{sec:results}.

% Key sources:
% \cite{Bowring2013_pjm}           -- RPM design and mitigation overview by IMM
% \cite{Hobbs2007_rpm}             -- Dynamic analysis of VRR / RPM
% \cite{Cramton2005_capacity}      -- Why capacity markets and demand curve design
% \cite{Joskow2007_capacity}       -- Reliability and competitive electricity markets
% \cite{Joskow2008_capacity_payments} -- Capacity payments need and design
% PJM Manual 18, PJM IMM SotM 2025 Q3
